\scanpage{29}
\mspage{27}
\deleted{Pren\deleted{\unsure{n}}ons en effet un schéma \inserted{$S$} modulaire à échelons \deleted{$S$} pour variations de courbes elliptiques, et localisons en $\underline{Z}_p$, soit encore $S$ le schéma lisse sur $\underline{Z}_p$ obtenu. On a sur $S$ le Module loc libre de rang 2 $H^1_{\operatorname{DR}}(A/S)$, et quand on enlève de $S$ le nombre fini de points fermés correspondants à Hasse nul, on aurait une stratification \deleted{absolue sur} rel. à $\underline{Z}_p$ sur ce Module. Pour des raisons de pureté évidentes, celle-ci se prolongerait au dessus de $S$ tout entier, d'où une stratification sur le $H^1_{\operatorname{DR}}$ de la famille modulaire de courbes elliptiques, en caractéristique $p$. Celle-ci serait alors compatible avec frobenius (car elle l'est dans le complémentaire de l'ensemble fini des points de Hasse nul), ce qui contredit nos réflexions de 2.7.}

Il n'est pas exclu entièrement, cependant, que en restant en car.$p>0$, l'énoncé sur la non-variation \deleted{\unsure{de $G(A)$}} infinitésimale de $G(A)$ avec $A$, $A$ partout ordinaire, soit vrai. Cela impliquerait que sur l'ouvert des valeurs ``ordinaires'' de l'invariant, le $H^1_{\operatorname{DR}}$ a une stratification canonique, mais celle-ci ne se prolongerait pas à la courbe modulaire toute entière. Mais je dois dire que ce drôle de comportement, où on aurait une stratification naturelle en car.$p$ et une autre en car. 0, sans qu'elles veuillent se recoller, semble assez canularesque.

\inserted{[}Les considérations précédentes font bien ressortir la nature ``infinie'' de la notion de stratification, par contraste avec celle de connexion, malgré les trompeuses apparences de la caractéristique nulle. Ainsi, sur le $H^1_{\operatorname{DR}}$ de la famille modulaire sur $\underline{Z}$ \deleted{\unsure{pade|par}} \inserted{de} schémas elliptiques, il y a \deleted{sur} \inserted{au dessus de} la fibre générique $S_{\underline{Q}}$ de $S$ une stratification naturelle, mais nous venons de voir (ou presque ...) que cette\deleted{b} stratification ne s'étend pas en une stratification sur $S_U$, $U$ un ouvert non vide de $\operatorname{Spec}(\underline{Z})$. (Le ``presque'' provient du fait qu'il n'est pas absolument clair si la stratification qu'on obtiendrait ainsi en car.$p$ serait respectée par frobenius; il faut absolument tirer au clair cet exemple particulier !). C'est en un sens assez moral, puisque pour nous la \inserted{module à} stratification doit jouer dans une large mesure le rôle d'un faisceau $\ell$-adique, dont il partage
\margin{28}
